\chapter*{Prólogo}
\addcontentsline{toc}{chapter}{Prólogo}
\markboth{PRÓLOGO}{}

Hay un momento, en casi toda tesis de pregrado que toca macroeconomía con
heterogeneidad, en que el asesor dice: ``esto está bien, pero es estático;
habría que hacerlo dinámico''. Y ahí se abre un abismo. El modelo estático se
resolvía con cuatro ecuaciones y un poco de álgebra. La versión dinámica, en
cambio, parece exigir dominar una literatura entera --- Krusell y Smith, Reiter,
proyección, perturbación, reducción de dimensionalidad --- antes de poder
escribir la primera línea de código.

Este manual existe porque ese abismo es más estrecho de lo que parece, siempre
que uno entre por la puerta correcta. La puerta correcta se llama \emph{jacobiano
en el espacio de secuencias}, y la idea que la sostiene cabe en un párrafo:

\begin{idea}
Todo lo que la heterogeneidad hace por el equilibrio general, a primer orden,
se resume en un conjunto de matrices. Cada matriz dice cuánto cambia una
variable agregada en la fecha $t$ cuando una variable agregada distinta cambia
en la fecha $s$. Una vez que uno tiene esas matrices, la distribución de agentes
--- que puede tener cientos de miles de puntos --- deja de aparecer en el
problema. El equilibrio general se vuelve álgebra lineal con matrices de tamaño
$T\times T$, donde $T$ es el horizonte, no el tamaño del espacio de estados.
\end{idea}

El aporte técnico de \ABRS\ fue encontrar una forma de calcular esas matrices
unas doscientas veces más rápido de lo que costaría hacerlo a la fuerza bruta.
A ese procedimiento lo llamaron \emph{algoritmo fake news}, por una razón que
quedará clara en el capítulo \ref{cap:fakenews} y que es, además, una de las
intuiciones más bonitas de la macroeconomía computacional reciente.

\paragraph{Qué supone este manual que el lector ya sabe}
Optimización dinámica a nivel de un curso intermedio: qué es una ecuación de
Bellman, qué es una ecuación de Euler, qué es un estado estacionario. Álgebra
lineal: matrices, inversas, autovalores, aunque los repasamos. Programación
básica en Python o en cualquier lenguaje con arreglos. \emph{No} supone haber
visto antes métodos numéricos para modelos de agentes heterogéneos; todo eso se
construye desde cero en el capítulo \ref{cap:bloque}.

\paragraph{Qué no es este manual}
No es un sustituto del artículo original, que es breve y está muy bien escrito;
es una rampa de acceso a él. Tampoco es un manual de programación: el código que
aparece está pensado para ser leído, no para competir en velocidad con
implementaciones profesionales. Y no cubre métodos globales ni perturbaciones de
orden superior, que es a donde hay que ir si uno necesita primas de riesgo,
elección de cartera o política óptima de Ramsey. En la sección
\ref{sec:limites-perturbacion} se explica por qué.

\paragraph{Una advertencia honesta sobre el modelo de bancos}
La parte V de este manual construye una extensión dinámica de un plan de trabajo
de tesis sobre encaje y heterogeneidad bancaria en el Perú. Esa extensión es un
\emph{prototipo pedagógico}: sirve para mostrar cómo se traduce un modelo
estático al espacio de secuencias y qué preguntas nuevas aparecen al hacerlo.
No es una calibración del sistema bancario peruano, y varios de sus supuestos
--- en particular la forma reducida de la fricción de pago de dividendos ---
tendrían que reemplazarse antes de usar sus números para decir algo cuantitativo.
Dónde exactamente, y con qué, está dicho en la sección \ref{sec:limitaciones-prototipo}.
